%%%%%%%%%%%%%%%%
%%% lev 11 %%%
%%%%%%%%%%%%%%%%

\Chapter{11}

\Verse{1}
{
	\hP{וַיְדַבֵּר יְהֹוָה אֶל מֹשֶׁה וְאֶל אַהֲרֹן לֵאמֹר אֲלֵהֶם׃}
}
{
	\eP{And YHWH spoke to Moses and to Aaron, saying to them,}
}
{
%	\footnote{
%		YHWH spoke to Moses and to Aaron, saying. This is only the second report
%		of laws being addressed to Aaron along with Moses. The first law given
%		to the people, the Passover command, is pronounced to both Moses and
%		Aaron (Exod 12:1), but then the laws are stated only to Moses. Indeed,
%		the only other time that YHWH has spoken to both Moses and Aaron until
%		now was at the beginning of their mission in Egypt (and there the text
%		does not report the words that God says to them). The inclusion of Aaron
%		at this point may be more of the divine response to Aaron’s loss of his
%		sons (see comment on 10:8). Or it may be because God has just told Aaron
%		that he and his sons are to instruct the people and distinguish between
%		holy and secular and between impure and clean (10:10–11), and so now God
%		begins to inform Aaron of the laws that he needs to know.
%	}
}

\Verse{2}
{
	\hP{דַּבְּרוּ אֶל בְּנֵי יִשְׂרָאֵל לֵאמֹר זֹאת הַחַיָּה אֲשֶׁר תֹּאכְלוּ מִכּל הַבְּהֵמָה אֲשֶׁר עַל הָאָרֶץ׃}
}
{
	\eP{
		“Speak to the children of Israel, saying: This is the
		animal that you shall eat out of all the domestic animals
		that are on the earth:
	}
}
{
%	\footnote{
%		This is the animal that you shall eat. The laws of permitted and
%		forbidden animals add further restrictions to the sacrificial laws. God
%		requires not only ritual slaughter, priestly performance, and location
%		at the Tent of Meeting; God delineates to Moses and Aaron which animals
%		may be eaten as well. Among land animals: those that have split hooves
%		and chew their cud. In water: those with scales and fins. (Redundant?
%		All scaled fish have fins.) In the air: the forbidden birds are named
%		one by one, and all flying creatures that have four legs are forbidden.
%		Also forbidden are creatures that have numerous legs, those that crawl
%		on their bellies, and any animals at all that have died on their own or
%		that have been killed by other animals. The reasons for these particular
%		inclusions and exclusions have never been worked out persuasively.
%		Explanations based on health and hygiene are difficult to defend, both
%		because the text never states this and because such explanations cannot
%		be consistently applied to the cases. Popular claims about the unhealthy
%		habits of pigs, for example, apply equally to various fowl that are
%		permitted. Some suggest that the Israelites may have observed
%		empirically that people who ate pork suffered more illness than others,
%		but this is without any textual or historical evidence; nor does it
%		explain how this fact was apparent to Israelites but unobserved by the
%		rest of the world for thousands of years. Popular claims that the
%		forbidden animals were objects of worship are simply wrong. Suggestions
%		that the system is ultimately about matters of ethics and not food—that
%		is, that it teaches discipline, self-denial, and so on—are not supported
%		in the text. The text rather relates various food laws to issues of
%		impurity (see below) and “abomination,” which fall in the realm of
%		ritual, not ethical, categories. Observance of ritual laws may have an
%		ethical dimension, such as discipline, but this may be said of all
%		ritual laws. It sheds no special light on this particular group of
%		commandments. The anthropologist Mary Douglas has suggested what is
%		perhaps the most-cited treatment at present (Purity and Danger), which
%		she and I have discussed: that the issue is one of defined classes of
%		animals: Split-hoofed, cud-chewing animals are “the model of the proper
%		kind of food for a pastoralist” (p. 54), and so this became the model of
%		what was appropriate to eat. Wild goats and antelope share these two
%		characteristics and therefore were acceptable. Pigs, which have a split
%		hoof but do not chew their cud, were excluded originally for the sole
%		reason that they did not fit this class as defined. Similarly with the
%		forbidden water creatures: they did not fit the class that had the
%		“right” kind of locomotion, namely fins and scales. Similarly with
%		animals that have two hands and two legs (weasels, mice, lizards) but
%		walk on all fours like a quadruped: they “perversely” use hands for
%		walking. Not having the right form of locomotion, they are forbidden.
%		Douglas’s explanation appears to me to be a self-defining system.
%		Whatever is permitted she identifies as what the Israelites perceived to
%		be “right.” Whatever is forbidden, she concludes, came to be forbidden
%		because it did not fit the category. This is mere definitional
%		description unless it can shed some light on the underlying reasons for
%		these perceptions of what was proper. Where Douglas tries to identify
%		the reasons she errs, in my judgment. The assertion that the split hoof
%		and the chewing of cud were “the model of the proper kind of food for a
%		pastoralist” is unhelpful without proof that Israel was in fact
%		primarily a pastoral society. It was not. (And Leviticus 11 certainly
%		does not appear to come from a pastoral circle. It comes from an urban
%		priesthood.) Nor does this explain why the non-Israelite residents of
%		the land, who alternately (and sometimes simultaneously) lived on the
%		same sites, conceived of these same animals so differently. Excavations
%		uncover pig bones in a Canaanite layer of a site, then no pig bones in
%		an Israelite layer of the same site, then pig bones again in the site’s
%		next Canaanite layer. There is likewise no reason to associate the
%		presence of scales on a fish with a concept of a proper form of
%		locomotion. And Douglas’s notion of a “perverse” form of locomotion by
%		animals whose front limbs are handlike is simply a misunderstanding of
%		the Hebrew term (kap), which does not mean “hand.” It refers elsewhere
%		to both the hands and the feet (Gen 8:9; Deut 2:5; 11:24; 28:35,56; Josh
%		1:3; 2 Sam 14:25) and is correctly translated “paws” (a translation that
%		Douglas appears to eschew). In its context here it follows a verse that
%		refers to the matter of split hooves (11:26), and it is simply
%		forbidding all animals that have soft paws rather than hooves. Douglas’s
%		treatment may add some helpful perspectives to the discussion, but the
%		essential mystery remains unsolved. I have tried here to exclude some of
%		the common explanations that appear to me to be looking in the wrong
%		directions, and I mean to identify some elements that should be taken
%		into account. The first is: similarity to humans. The fact that the text
%		elsewhere has conceived of animals, but not plants, as having the
%		quality of life (nepeš) in common with humans, even though plants do in
%		fact live and die, already means that the degree of similarity to humans
%		figured in the distinctions, whether consciously or not. Plants were so
%		different from humans as not to be thought of as a life-form in the same
%		sense and therefore not even considered as having to be permitted or
%		forbidden. In the animal kingdom, meanwhile, on a continuum of likeness
%		to humans, at one extreme would be cannibalism, which is so extreme as
%		not even to be considered (unless the law against human sacrifice is
%		understood necessarily to prevent it). Next would be the ape family,
%		which is forbidden, having soft hands and feet; then all fellow
%		creatures having soft hands and feet or paws (kap) like humans rather
%		than hooves. Next might be all animals that, like humans, are
%		carnivorous (thus including numerous land animals, all birds of prey,
%		and various reptiles). The permissible animals would then be those that
%		are twice-distinguished from humans, both having a split hoof and
%		chewing their cud. Animals that have only one of these are especially
%		singled out by species, such as pigs (which have a split hoof but do not
%		chew the cud) and camels (which chew the cud but do not have a split
%		hoof). This does not explain all of the permissions and prohibitions of
%		Leviticus 11. Notably, it does not account for water creatures (few
%		creatures are as unlike humans as lobsters) or insects (why are locusts
%		permitted?). It does, however, relate to enough of the distinctions that
%		it should be taken into account in any attempt to explain this list.
%		Ultimately, since no single underlying principle has been discovered
%		that accounts for all the distinctions, it appears likely that there is
%		a convergence of two or probably more factors. It could be a combination
%		of a principle (such as likeness to humans) and completely idiosyncratic
%		factors (a distaste, phobia, or even allergy on the part of some
%		individual in an authoritative position). For the larger purpose of
%		coming to an understanding of Leviticus, it is important for now to note
%		the fact of distinction and how seriously and pervasively it is
%		developed here in regard to a basic function of life. The list confirms
%		this explicitly by concluding: This is the instruction [tôrh] for animal
%		and bird and every living being that moves in the water and for every
%		being that swarms on the earth, to distinguish[lhabdîl] between the
%		impure and the pure, and between the living thing that is eaten and the
%		living thing that shall not be eaten. (11:46–47)
%	}
}

\Verse{3}
{
	\hP{כֹּל מַפְרֶסֶת פַּרְסָה וְשֹׁסַעַת שֶׁסַע פְּרָסֹת מַעֲלַת גֵּרָה בַּבְּהֵמָה אֹתָהּ תֹּאכֵלוּ׃}
}
{
	\eP{
		every one that has a hoof and that has a split of hooves,
		that regurgitates cud among animals, you shall eat it.
	}
}
{
%	\footnote{
%		a split of hooves. Each of its hooves is split, e.g., cows; as opposed
%		to those that have a single hoof, e.g., horses.
%	}
}

\Verse{4}
{
	\hP{אַךְ אֶת זֶה לֹא תֹאכְלוּ מִמַּעֲלֵי הַגֵּרָה וּמִמַּפְרִסֵי הַפַּרְסָה אֶת הַגָּמָל כִּי מַעֲלֵה גֵרָה הוּא וּפַרְסָה אֵינֶנּוּ מַפְרִיס טָמֵא הוּא לָכֶם׃}
}
{
	\eP{
		Except you shall not eat this out of those that
		regurgitate the cud and out of those that have a hoof: The
		camel, because it regurgitates cud and does not have a
		hoof; it is impure to you.
	}
}
{
}

\Verse{5}
{
	\hP{וְאֶת הַשָּׁפָן כִּי מַעֲלֵה גֵרָה הוּא וּפַרְסָה לֹא יַפְרִיס טָמֵא הוּא לָכֶם׃}
}
{
	\eP{
		And the rock-badger, because it regurgitates cud and does
		not have a hoof; it is impure to you.
	}
}
{
}

\Verse{6}
{
	\hP{וְאֶת הָאַרְנֶבֶת כִּי מַעֲלַת גֵּרָה הִוא וּפַרְסָה לֹא הִפְרִיסָה טְמֵאָה הִוא לָכֶם׃}
}
{
	\eP{
		And the hare, because it regurgitates cud and does not
		have a hoof; it is impure to you.
	}
}
{
%	\footnote{
%		the hare, because it regurgitates cud. Hares do not regurgitate cud.
%	}
}

\Verse{7}
{
	\hP{וְאֶת הַחֲזִיר כִּי מַפְרִיס פַּרְסָה הוּא וְשֹׁסַע שֶׁסַע פַּרְסָה וְהוּא גֵּרָה לֹא יִגָּר טָמֵא הוּא לָכֶם׃}
}
{
	\eP{
		And the pig, because it has a hoof and has a split of
		hooves, and it does not chew cud; it is impure to you.
	}
}
{
}

\Verse{8}
{
	\hP{מִבְּשָׂרָם לֹא תֹאכֵלוּ וּבְנִבְלָתָם לֹא תִגָּעוּ טְמֵאִים הֵם לָכֶם׃}
}
{
	\eP{
		You shall not eat from their meat, and you shall not touch
		their carcass. They are impure to you.
	}
}
{
}

\Verse{9}
{
	\hP{אֶת זֶה תֹּאכְלוּ מִכֹּל אֲשֶׁר בַּמָּיִם כֹּל אֲשֶׁר לוֹ סְנַפִּיר וְקַשְׂקֶשֶׂת בַּמַּיִם בַּיַּמִּים וּבַנְּחָלִים אֹתָם תֹּאכֵלוּ׃}
}
{
	\eP{
		“You shall eat this out of all that are in the water:
		every one that has fins and scales in the water—in the
		seas and in the streams—you shall eat them.
	}
}
{
%	\footnote{
%		all that are in the water. There is no differentiation of species of
%		water creatures in the entire Tanak. All fish are simply called dg. My
%		colleague Jacob Milgrom has proposed that the reason for this is that
%		ancient Israelites knew little of fish because the eastern Mediterranean
%		was not suitable for fish until the opening of the Suez Canal in our
%		era, which changed the ecology of that sea. But even if this were
%		so—that there were no fish in the eastern Mediterranean—the Israelites
%		still had the Red Sea, the Sea of Galilee, and the Jordan River. Surely
%		those who lived near these bodies of water knew one fish from another. A
%		more likely explanation is that the list of permitted and forbidden
%		animals came from the Jerusalem priesthood, and before the invention of
%		refrigeration it would have been extremely rare for anyone in Jerusalem
%		to have an opportunity to eat fish.
%	}
}

\Verse{10}
{
	\hP{וְכֹל אֲשֶׁר אֵין לוֹ סְנַפִּיר וְקַשְׂקֶשֶׂת בַּיַּמִּים וּבַנְּחָלִים מִכֹּל שֶׁרֶץ הַמַּיִם וּמִכֹּל נֶפֶשׁ הַחַיָּה אֲשֶׁר בַּמָּיִם שֶׁקֶץ הֵם לָכֶם׃}
}
{
	\eP{
		And any one that does not have fins and scales in the seas
		and in the streams—out of every swarming thing of the
		water and out of every living being that is in the
		water—they are a detestable thing to you,
	}
}
{
}

\Verse{11}
{
	\hP{וְשֶׁקֶץ יִהְיוּ לָכֶם מִבְּשָׂרָם לֹא תֹאכֵלוּ וְאֶת נִבְלָתָם תְּשַׁקֵּצוּ׃}
}
{
	\eP{
		and they shall be a detestable thing to you. You shall not
		eat from their meat, and you shall detest their carcass.
	}
}
{
}

\Verse{12}
{
	\hP{כֹּל אֲשֶׁר אֵין לוֹ סְנַפִּיר וְקַשְׂקֶשֶׂת בַּמָּיִם שֶׁקֶץ הוּא לָכֶם׃}
}
{
	\eP{
		Every one that does not have scales and fins in the water,
		it is a detestable thing to you.
	}
}
{
}

\Verse{13}
{
	\hP{וְאֶת אֵלֶּה תְּשַׁקְּצוּ מִן הָעוֹף לֹא יֵאָכְלוּ שֶׁקֶץ הֵם אֶת הַנֶּשֶׁר וְאֶת הַפֶּרֶס וְאֵת הָעזְנִיָּה׃}
}
{
	\eP{
		“And you shall detest these out of the flying creatures.
		They shall not be eaten. They are a detestable thing: the
		eagle and the vulture and the black vulture
	}
}
{
}

\Verse{14}
{
	\hP{וְאֶת הַדָּאָה וְאֶת הָאַיָּה לְמִינָהּ׃}
}
{
	\eP{and the kite and the falcon by its kind,}
}
{
}

\Verse{15}
{
	\hP{אֵת כּל עֹרֵב לְמִינוֹ׃}
}
{
	\eP{every raven by its kind}
}
{
}

\Verse{16}
{
	\hP{וְאֵת בַּת הַיַּעֲנָה וְאֶת הַתַּחְמָס וְאֶת הַשָּׁחַף וְאֶת הַנֵּץ לְמִינֵהוּ׃}
}
{
	\eP{
		and the eagle owl and the nighthawk and the seagull and
		the hawk by its kind
	}
}
{
%	\footnote{
%		eagle owl. Hebrew bat hayya’nh. It has usually been identified as an
%		ostrich. Milgrom has thrown that identification into doubt. I agree.
%		Moreover, ostriches are primarily plant-eating, like the permitted
%		birds. Therefore: the ostrich is not necessarily forbidden for food.
%	}
}

\Verse{17}
{
	\hP{וְאֶת הַכּוֹס וְאֶת הַשָּׁלָךְ וְאֶת הַיַּנְשׁוּף׃}
}
{
	\eP{and the little owl and the cormorant and the great owl}
}
{
}

\Verse{18}
{
	\hP{וְאֶת הַתִּנְשֶׁמֶת וְאֶת הַקָּאָת וְאֶת הָרָחָם׃}
}
{
	\eP{and the white owl and the pelican and the fish hawk}
}
{
}

\Verse{19}
{
	\hP{וְאֵת הַחֲסִידָה הָאֲנָפָה לְמִינָהּ וְאֶת הַדּוּכִיפַת וְאֶת הָעֲטַלֵּף׃}
}
{
	\eP{
		and the stork and the heron by its kind and the hoopoe and
		the bat.
	}
}
{
}

\Verse{20}
{
	\hP{כֹּל שֶׁרֶץ הָעוֹף הַהֹלֵךְ עַל אַרְבַּע שֶׁקֶץ הוּא לָכֶם׃}
}
{
	\eP{
		Every swarming thing of flying creatures that goes on
		four—it is a detestable thing to you.
	}
}
{
}

\Verse{21}
{
	\hP{אַךְ אֶת זֶה תֹּאכְלוּ מִכֹּל שֶׁרֶץ הָעוֹף הַהֹלֵךְ עַל אַרְבַּע אֲשֶׁר [לוֹ] (לא) כְרָעַיִם מִמַּעַל לְרַגְלָיו לְנַתֵּר בָּהֵן עַל הָאָרֶץ׃}
}
{
	\eP{
		Except you shall eat this out of every swarming thing of
		flying creatures that goes on four: that which has jointed
		legs above its feet, to leap by them on the earth.
	}
}
{
	Qere Ketiv.
}

\Verse{22}
{
	\hP{אֶת אֵלֶּה מֵהֶם תֹּאכֵלוּ אֶת הָאַרְבֶּה לְמִינוֹ וְאֶת הַסּלְעָם לְמִינֵהוּ וְאֶת הַחַרְגֹּל לְמִינֵהוּ וְאֶת הֶחָגָב לְמִינֵהוּ׃}
}
{
	\eP{
		You shall eat these out of them: the locust by its kind
		and the bald locust by its kind and the cricket by its
		kind and the grasshopper by its kind.
	}
}
{
}

\Verse{23}
{
	\hP{וְכֹל שֶׁרֶץ הָעוֹף אֲשֶׁר לוֹ אַרְבַּע רַגְלָיִם שֶׁקֶץ הוּא לָכֶם׃}
}
{
	\eP{
		And every swarming thing of flying creatures that has four
		legs—it is a detestable thing to you.
	}
}
{
}

\Verse{24}
{
	\hP{וּלְאֵלֶּה תִּטַּמָּאוּ כּל הַנֹּגֵעַ בְּנִבְלָתָם יִטְמָא עַד הָעָרֶב׃}
}
{
	\eP{
		“And you will become impure by these—everyone who touches
		their carcass will be impure until evening,
	}
}
{
}

\Verse{25}
{
	\hP{וְכל הַנֹּשֵׂא מִנִּבְלָתָם יְכַבֵּס בְּגָדָיו וְטָמֵא עַד הָעָרֶב׃}
}
{
	\eP{
		and everyone who carries their carcass shall wash his
		clothes and be impure until evening—
	}
}
{
}

\Verse{26}
{
	\hP{לְכל הַבְּהֵמָה אֲשֶׁר הִוא מַפְרֶסֶת פַּרְסָה וְשֶׁסַע אֵינֶנָּה שֹׁסַעַת וְגֵרָה אֵינֶנָּה מַעֲלָה טְמֵאִים הֵם לָכֶם כּל הַנֹּגֵעַ בָּהֶם יִטְמָא׃}
}
{
	\eP{
		by every animal that has a hoof and does not have a split
		of hooves and does not regurgitate its cud. They are
		impure to you. Everyone who touches them will be impure.
	}
}
{
}

\Verse{27}
{
	\hP{וְכֹל הוֹלֵךְ עַל כַּפָּיו בְּכל הַחַיָּה הַהֹלֶכֶת עַל אַרְבַּע טְמֵאִים הֵם לָכֶם כּל הַנֹּגֵעַ בְּנִבְלָתָם יִטְמָא עַד הָעָרֶב׃}
}
{
	\eP{
		And every one that goes on its paws among all the animals
		that go on four, they are impure to you. Everyone who
		touches their carcass will be impure until evening,
	}
}
{
}

\Verse{28}
{
	\hP{וְהַנֹּשֵׂא אֶת נִבְלָתָם יְכַבֵּס בְּגָדָיו וְטָמֵא עַד הָעָרֶב טְמֵאִים הֵמָּה לָכֶם׃}
}
{
	\eP{
		and one who carries their carcass shall wash his clothes
		and be impure until evening. They are impure to you.
	}
}
{
}

\Verse{29}
{
	\hP{וְזֶה לָכֶם הַטָּמֵא בַּשֶּׁרֶץ הַשֹּׁרֵץ עַל הָאָרֶץ הַחֹלֶד וְהָעַכְבָּר וְהַצָּב לְמִינֵהוּ׃}
}
{
	\eP{
		“And this is the impure to you among the things that swarm
		on the earth: the rat, the mouse, and the great lizard by
		its kind
	}
}
{
}

\Verse{30}
{
	\hP{וְהָאֲנָקָה וְהַכֹּחַ וְהַלְּטָאָה וְהַחֹמֶט וְהַתִּנְשָׁמֶת׃}
}
{
	\eP{
		and the gecko and the spotted lizard and the lizard and
		the sand lizard and the chameleon.
	}
}
{
%	\footnote{
%		spotted lizard. The identification of several of the animals on this
%		list is uncertain. Milgrom favors the spotted lizard for this one.
%	}
}

\Verse{31}
{
	\hP{אֵלֶּה הַטְּמֵאִים לָכֶם בְּכל הַשָּׁרֶץ כּל הַנֹּגֵעַ בָּהֶם בְּמֹתָם יִטְמָא עַד הָעָרֶב׃}
}
{
	\eP{
		These are the impure to you among all the swarming things.
		Everyone who touches them when they are dead will be
		impure until evening.
	}
}
{
}

\Verse{32}
{
	\hP{וְכֹל אֲשֶׁר יִפֹּל עָלָיו מֵהֶם בְּמֹתָם יִטְמָא מִכּל כְּלִי עֵץ אוֹ בֶגֶד אוֹ עוֹר אוֹ שָׂק כּל כְּלִי אֲשֶׁר יֵעָשֶׂה מְלָאכָה בָּהֶם בַּמַּיִם יוּבָא וְטָמֵא עַד הָעֶרֶב וְטָהֵר׃}
}
{
	\eP{
		And everything on which one of them will fall when they
		are dead will be impure, out of every wooden item or
		clothing or leather or sack, every item with which work
		may be done shall be put in water and will be impure until
		evening and then will be pure.
	}
}
{
}

\Verse{33}
{
	\hP{וְכל כְּלִי חֶרֶשׂ אֲשֶׁר יִפֹּל מֵהֶם אֶל תּוֹכוֹ כֹּל אֲשֶׁר בְּתוֹכוֹ יִטְמָא וְאֹתוֹ תִשְׁבֹּרוּ׃}
}
{
	\eP{
		And every clay container into which any of them will fall:
		everything that is inside it will be impure, and you shall
		break it.
	}
}
{
}

\Verse{34}
{
	\hP{מִכּל הָאֹכֶל אֲשֶׁר יֵאָכֵל אֲשֶׁר יָבוֹא עָלָיו מַיִם יִטְמָא וְכל מַשְׁקֶה אֲשֶׁר יִשָּׁתֶה בְּכל כְּלִי יִטְמָא׃}
}
{
	\eP{
		Every food that may be eaten on which water will come will
		be impure, and every liquid that may be drunk in any
		container will be impure.
	}
}
{
}

\Verse{35}
{
	\hP{וְכֹל אֲשֶׁר יִפֹּל מִנִּבְלָתָם עָלָיו יִטְמָא תַּנּוּר וְכִירַיִם יֻתָּץ טְמֵאִים הֵם וּטְמֵאִים יִהְיוּ לָכֶם׃}
}
{
	\eP{
		And everything on which some of their carcass will fall
		will be impure. An oven or stove shall be demolished. They
		are impure, and they shall be impure to you.
	}
}
{
}

\Verse{36}
{
	\hP{אַךְ מַעְיָן וּבוֹר מִקְוֵה מַיִם יִהְיֶה טָהוֹר וְנֹגֵעַ בְּנִבְלָתָם יִטְמָא׃}
}
{
	\eP{
		Except a spring or a cistern with a concentration of
		waters will be pure. And one who touches their carcass
		will be impure.
	}
}
{
}

\Verse{37}
{
	\hP{וְכִי יִפֹּל מִנִּבְלָתָם עַל כּל זֶרַע זֵרוּעַ אֲשֶׁר יִזָּרֵעַ טָהוֹר הוּא׃}
}
{
	\eP{
		And if some of their carcass will fall on any sowing seed
		that will be sown, it is pure.
	}
}
{
}

\Verse{38}
{
	\hP{וְכִי יֻתַּן מַיִם עַל זֶרַע וְנָפַל מִנִּבְלָתָם עָלָיו טָמֵא הוּא לָכֶם׃}
}
{
	\eP{
		And if water will be put on a seed, and some of their
		carcass will fall on it, it is impure to you.
	}
}
{
}

\Verse{39}
{
	\hP{וְכִי יָמוּת מִן הַבְּהֵמָה אֲשֶׁר הִיא לָכֶם לְאכְלָה הַנֹּגֵעַ בְּנִבְלָתָהּ יִטְמָא עַד הָעָרֶב׃}
}
{
	\eP{
		“And if any of the animals that are yours for food will
		die, the one who touches its carcass will be impure until
		evening,
	}
}
{
}

\Verse{40}
{
	\hP{וְהָאֹכֵל מִנִּבְלָתָהּ יְכַבֵּס בְּגָדָיו וְטָמֵא עַד הָעָרֶב וְהַנֹּשֵׂא אֶת נִבְלָתָהּ יְכַבֵּס בְּגָדָיו וְטָמֵא עַד הָעָרֶב׃}
}
{
	\eP{
		and the one who eats some of its carcass shall wash his
		clothes and is impure until evening, and the one who
		carries its carcass shall wash his clothes and is impure
		until evening.
	}
}
{
}

\Verse{41}
{
	\hP{וְכל הַשֶּׁרֶץ הַשֹּׁרֵץ עַל הָאָרֶץ שֶׁקֶץ הוּא לֹא יֵאָכֵל׃}
}
{
	\eP{
		“And every swarming creature that swarms on the earth, it
		is a detestable thing. It shall not be eaten.
	}
}
{
}

\Verse{42}
{
	\hP{כֹּל הוֹלֵךְ עַל גָּחוֹן וְכֹל הוֹלֵךְ עַל אַרְבַּע עַד כּל מַרְבֵּה רַגְלַיִם לְכל הַשֶּׁרֶץ הַשֹּׁרֵץ עַל הָאָרֶץ לֹא תֹאכְלוּם כִּי שֶׁקֶץ הֵם׃}
}
{
	\eP{
		Everything going on a belly and everything going on four
		as well as everything having a great number of legs of
		every swarming creature that swarms on the earth: you
		shall not eat them because they are a detestable thing.
	}
}
{
%	\footnote{
%		belly. The third letter of this word in the Hebrew text is written extra
%		large, thus: —and it is the middle of the Torah in terms of letters.
%	}
}

\Verse{43}
{
	\hP{אַל תְּשַׁקְּצוּ אֶת נַפְשֹׁתֵיכֶם בְּכל הַשֶּׁרֶץ הַשֹּׁרֵץ וְלֹא תִטַּמְּאוּ בָּהֶם וְנִטְמֵתֶם בָּם׃}
}
{
	\eP{
		Do not make yourselves detestable with any swarming
		creature that swarms, and you shall not become impure with
		them, so that you will be impure through them.
	}
}
{
}

\Verse{44}
{
	\hP{כִּי אֲנִי יְהֹוָה אֱלֹהֵיכֶם וְהִתְקַדִּשְׁתֶּם וִהְיִיתֶם קְדֹשִׁים כִּי קָדוֹשׁ אָנִי וְלֹא תְטַמְּאוּ אֶת נַפְשֹׁתֵיכֶם בְּכל הַשֶּׁרֶץ הָרֹמֵשׂ עַל הָאָרֶץ׃}
}
{
	\eP{
		Because I am YHWH, your {\God}, and you shall make yourselves
		holy, and you shall be holy because I am holy. And you
		shall not make yourselves impure with any swarming
		creature that creeps on the earth.
	}
}
{
}

\Verse{45}
{
	\hP{כִּי אֲנִי יְהֹוָה הַמַּעֲלֶה אֶתְכֶם מֵאֶרֶץ מִצְרַיִם לִהְיֹת לָכֶם לֵאלֹהִים וִהְיִיתֶם קְדֹשִׁים כִּי קָדוֹשׁ אָנִי׃}
}
{
	\eP{
		Because I am YHWH, who brought you up from the land of
		Egypt to be {\God} to you, and you shall be holy because I am
		holy.
	}
}
{
}

\Verse{46}
{
	\hP{זֹאת תּוֹרַת הַבְּהֵמָה וְהָעוֹף וְכֹל נֶפֶשׁ הַחַיָּה הָרֹמֶשֶׂת בַּמָּיִם וּלְכל נֶפֶשׁ הַשֹּׁרֶצֶת עַל הָאָרֶץ׃}
}
{
	\eP{
		“This is the instruction for animal and bird and every
		living being that moves in the water and for every being
		that swarms on the earth,
	}
}
{
}

\Verse{47}
{
	\hP{לְהַבְדִּיל בֵּין הַטָּמֵא וּבֵין הַטָּהֹר וּבֵין הַחַיָּה הַנֶּאֱכֶלֶת וּבֵין הַחַיָּה אֲשֶׁר לֹא תֵאָכֵל׃}
}
{
	\eP{
		to distinguish between the impure and the pure, and
		between the living thing that is eaten and the living
		thing that shall not be eaten.”
	}
}
{
}
